% ========================================================================================
% CAPITOLO 3 - MACCHINE ARITMETICHE
% ========================================================================================
\chapter{Macchine aritmetiche}

% ----------------------------------------------------------------------------------------
\section{Esercizio 7: Moltiplicatore di Booth}
% ----------------------------------------------------------------------------------------

\subsection{Esercizio 7.1}

\begin{traccia}
	Progettare, implementare in VHDL e simulare una macchina moltiplicatore di Booth in grado
	di effettuare il prodotto di 2 stringhe A e B da 8 bit ciascuna.
\end{traccia}

\progetto

L'algoritmo di Booth (radix-2) moltiplica due interi in complemento a 2 esaminando a ogni
passo la coppia di bit $(Q_0, Q_{-1})$ del moltiplicatore (Algoritmo~\ref{alg:booth}).

\begin{algorithm}[H]
	\caption{Moltiplicazione di Booth radix-2 su $n$ bit.}
	\label{alg:booth}
	\KwIn{moltiplicando $M$ e moltiplicatore $Q$, in complemento a 2 su $n$ bit}
	\KwOut{prodotto $P$ su $2n$ bit}
	$A \leftarrow 0$;\quad $Q_{-1} \leftarrow 0$;\quad $count \leftarrow n$\;
	\While{$count > 0$}{
		\uIf{$Q_0 Q_{-1} = 10$}{
			$A \leftarrow A - M$\;
		}
		\ElseIf{$Q_0 Q_{-1} = 01$}{
			$A \leftarrow A + M$\;
		}
		shift aritmetico a destra di $(A, Q, Q_{-1})$\;
		$count \leftarrow count - 1$\;
	}
	$P \leftarrow (A, Q)$\;
\end{algorithm}

\begin{attenzione}[Il caso $M = -128$]
	Con $A$ su 8 bit, l'operazione $A - M$ con $M = -128$ esce dal range rappresentabile.
	\todo{Mostrare come è gestito: accumulatore su $n+1$ bit oppure dimostrazione che il caso non si presenta.}
\end{attenzione}

\todo{Esempio svolto a mano, schema PO/PC, ASM chart dell'unità di controllo.}

\implementazione
\todo{Codice di parte operativa e unità di controllo.}

\simulazione
\todo{Testbench esaustivo sulle 65\,536 coppie, confronto con \texttt{signed(A) * signed(B)}.}

\timinganalysis
\todo{Percorso critico attraverso il sommatore/sottrattore, frequenza massima.}

\subsection{Esercizio 7.2}

\begin{traccia}
	Sintetizzare il moltiplicatore implementato al punto 7.1 su FPGA e testarlo mediante
	l'utilizzo dei dispositivi di input/output (switch, bottoni, led, display) presenti sulla
	board di sviluppo in dotazione. La modalità di utilizzo degli stessi è a completa
	discrezione degli studenti.
\end{traccia}

\sintesiboard
\todo{Operandi sugli switch, start su pulsante, prodotto sul display.}

% ----------------------------------------------------------------------------------------
\section{Esercizio 7BIS: Divisore Non-Restoring (solo 9 CFU)}
% ----------------------------------------------------------------------------------------

\subsection{Esercizio 7BIS.1}

\begin{traccia}
	Progettare, implementare in VHDL e simulare una macchina divisore\linebreak[0] (modalità non-restoring)
	in grado di effettuare la divisione intera fra due stringhe A e B di 4 bit ciascuna.
\end{traccia}

\progetto
\todo{Algoritmo non-restoring, correzione finale del resto, divisione per zero.}

\implementazione
\todo{Codice.}

\simulazione
\todo{Testbench esaustivo confrontato con \texttt{/} e \texttt{mod}.}

\subsection{Esercizio 7BIS.2}

\begin{traccia}
	Sintetizzare il divisore implementato al punto 7BIS.1 su FPGA e testarlo mediante
	l'utilizzo dei dispositivi di input/output (switch, bottoni, led, display) presenti sulla
	board di sviluppo in dotazione. La modalità di utilizzo degli stessi è a completa
	discrezione degli studenti.
\end{traccia}

\sintesiboard
\todo{Top-level, vincoli, prova su scheda.}
